\section{Praktik: Merancang dan Mengelola Basis Data untuk Aplikasi Web}

Merancang database yang baik memerlukan tahapan: (1) identifikasi \textbf{entitas} (objek utama: pengguna, produk, pesanan); (2) tentukan \textbf{atribut} (kolom) setiap entitas; (3) identifikasi \textbf{relasi} antar entitas dan tentukan kardinalitas (one-to-one, one-to-many, many-to-many); (4) terapkan \textbf{normalisasi} untuk mengurangi redundansi; (5) buat \textbf{diagram ER} sebelum menulis SQL; (6) tentukan \textbf{index} pada kolom yang sering dicari; (7) atur \textbf{constraint} (NOT NULL, UNIQUE, FK) untuk integritas data \cite{mysqlDocs}.

\begin{figure}[ht]
\centering
\begin{tikzpicture}[
  entity/.style={draw, rectangle, rounded corners, minimum width=2.5cm, minimum height=1.3cm, align=center, font=\footnotesize, fill=#1},
  arrow/.style={thick}
]
  \node[entity=blue!15] (user) {\textbf{pengguna}\\id (PK), username,\\email, password};
  \node[entity=green!15, right=2cm of user] (order) {\textbf{pesanan}\\id (PK), user\_id (FK),\\tanggal, total, status};
  \node[entity=orange!15, right=2cm of order] (detail) {\textbf{pesanan\_detail}\\id (PK), pesanan\_id (FK),\\produk\_id (FK), jumlah};
  \node[entity=red!10, below=1.5cm of detail] (produk) {\textbf{produk}\\id (PK), nama, harga,\\kategori\_id (FK)};
  \node[entity=purple!10, below=1.5cm of user] (kat) {\textbf{kategori}\\id (PK), nama};

  \draw[arrow] (user) -- node[above, font=\scriptsize]{1:N} (order);
  \draw[arrow] (order) -- node[above, font=\scriptsize]{1:N} (detail);
  \draw[arrow] (detail) -- node[right, font=\scriptsize]{N:1} (produk);
  \draw[arrow] (produk) -- node[below, font=\scriptsize]{N:1} (kat);
\end{tikzpicture}
\caption{Diagram ER untuk aplikasi e-commerce sederhana}
\end{figure}

\begin{webcode}[language=SQL, caption={SQL lengkap untuk skema e-commerce}]
CREATE DATABASE IF NOT EXISTS ecommerce CHARACTER SET utf8mb4;
USE ecommerce;

CREATE TABLE kategori (
  id   INT AUTO_INCREMENT PRIMARY KEY,
  nama VARCHAR(100) NOT NULL UNIQUE
);

CREATE TABLE produk (
  id          INT AUTO_INCREMENT PRIMARY KEY,
  nama        VARCHAR(200) NOT NULL,
  harga       DECIMAL(12,2) NOT NULL,
  stok        INT NOT NULL DEFAULT 0,
  kategori_id INT,
  FOREIGN KEY (kategori_id) REFERENCES kategori(id)
);

CREATE TABLE pengguna (
  id       INT AUTO_INCREMENT PRIMARY KEY,
  username VARCHAR(50) NOT NULL UNIQUE,
  email    VARCHAR(100) NOT NULL UNIQUE,
  password VARCHAR(255) NOT NULL
);

CREATE TABLE pesanan (
  id      INT AUTO_INCREMENT PRIMARY KEY,
  user_id INT NOT NULL,
  tanggal DATETIME DEFAULT CURRENT_TIMESTAMP,
  total   DECIMAL(12,2) DEFAULT 0,
  status  ENUM('pending','proses','selesai','batal') DEFAULT 'pending',
  FOREIGN KEY (user_id) REFERENCES pengguna(id)
);

CREATE TABLE pesanan_detail (
  id         INT AUTO_INCREMENT PRIMARY KEY,
  pesanan_id INT NOT NULL,
  produk_id  INT NOT NULL,
  jumlah     INT NOT NULL DEFAULT 1,
  harga      DECIMAL(12,2) NOT NULL,
  FOREIGN KEY (pesanan_id) REFERENCES pesanan(id) ON DELETE CASCADE,
  FOREIGN KEY (produk_id) REFERENCES produk(id)
);

CREATE INDEX idx_produk_kategori ON produk(kategori_id);
CREATE INDEX idx_pesanan_user ON pesanan(user_id);
\end{webcode}

